% Template for Matematika Dasar practice files.
% Copy this file into a topic folder and rename it as Set N.tex or Kuis N.tex.
% Leave \kuisfalse for open-ended Sets. Change to \kuistrue for multiple-choice Kuis.

\documentclass[12pt,a4paper]{article}

\usepackage[utf8]{inputenc}
\usepackage[T1]{fontenc}
\usepackage{amsmath}
\usepackage{amssymb}
\usepackage{geometry}
\usepackage{enumitem}

\geometry{margin=2.4cm}
\setlength{\parindent}{0pt}
\setlength{\parskip}{6pt}

\newif\ifkuis
\kuisfalse

\newcommand{\nomor}{N}
\newcommand{\mapel}{Matematika Dasar}
\newcommand{\topik}{Problem Topic}
\newcommand{\tanggal}{Hari - DD/MM/YYYY}
\newcommand{\setlevel}{easy / intermediate / difficult}

\ifkuis
	\newcommand{\jenis}{Kuis}
	\newcommand{\metalevel}{Tipe: Pilihan ganda, campuran intermediate--difficult}
\else
	\newcommand{\jenis}{Set}
	\newcommand{\metalevel}{Tingkat: \setlevel}
\fi

\newenvironment{soal}
	{\begin{enumerate}[label=\arabic*., leftmargin=*, itemsep=1.1em, topsep=0.5em]}
	{\end{enumerate}}

\newenvironment{pilihan}
	{\begin{enumerate}[label=\Alph*., leftmargin=2.5em, itemsep=0.2em, topsep=0.3em]}
	{\end{enumerate}}

\newcommand{\sheetheader}{
	\begin{center}
		{\Large\textbf{\jenis\ \nomor\ - \mapel}}\\[4pt]
		{\large \topik}\\[4pt]
		Tanggal: \tanggal
	\end{center}
	\vspace{0.4em}
	\hrule
	\vspace{0.8em}
	\textbf{\metalevel}
	\vspace{0.6em}
}

\begin{document}

\sheetheader

\ifkuis

\begin{soal}
	\item \textit{[Problem]}
	\begin{pilihan}
		\item \textit{[Option A]}
		\item \textit{[Option B]}
		\item \textit{[Option C]}
		\item \textit{[Option D]}
		\item \textit{[Option E]}
	\end{pilihan}

	\item \textit{[Problem]}
	\begin{pilihan}
		\item \textit{[Option A]}
		\item \textit{[Option B]}
		\item \textit{[Option C]}
		\item \textit{[Option D]}
		\item \textit{[Option E]}
	\end{pilihan}
\end{soal}

\else

\begin{soal}
	\item \textit{[Open-ended problem]}

	\item \textit{[Open-ended problem]}

	\item \textit{[Open-ended problem]}
\end{soal}

\fi

\end{document}
